\section{Objective}\label{sec:objective}

\sumpara The planner's primary measure is total weighted tardiness at the last stage. \sumparaend
\fin A job's tardiness is the time by which its completion at the last stage exceeds its due date, and it is zero when the job finishes on time. \fine
\fin The primary measure is total weighted tardiness, the sum over jobs of weight times tardiness, in weighted minutes. \fine
\fin Due dates are soft targets, so a late job makes a schedule worse without making it infeasible. \fine

\sumpara The secondary measure is the waiting that builds up between stages. \sumparaend
\fin Total waiting sums, over jobs and over every stage after the first, the time from a job's completion at one stage to its start at the next. \fine
\fin The course project allows unlimited waiting space between stages, so waiting measures work in process piling up on the line and carries no feasibility consequence. \fine
\fin Every comparison also reports the computation time of each method. \fine

\sumpara The methods optimize different things, and every reported score is recomputed after the fact. \sumparaend
\fin Iterated greedy compares sequences by total weighted tardiness and breaks exact ties by total waiting, because many sequences can share a tardiness value when most jobs finish on time. \fine
\fin The earliest due date rule optimizes nothing; it sorts the jobs once. \fine
\fin The mixed-integer program minimizes total weighted tardiness alone. \fine
\fin Every reported value of either measure is recomputed after the fact from the exported schedule, so no method reports its own score. \fine
